\chapter{LLM agents and tool use}
\label{ch:agents}

\begin{quote}
\emph{``An agent is an LLM that can take actions, not just generate text. It reasons about what to do, uses tools, and iterates until the task is complete.''}
\end{quote}

% ============================================================
\section{What is an LLM agent?}

A standard LLM takes a prompt and returns text. An \textbf{agent} adds:
\begin{itemize}
  \item \textbf{Tools}: functions the LLM can call (search, calculator, code execution, APIs).
  \item \textbf{Reasoning}: the LLM decides which tool to use and in what order.
  \item \textbf{Memory}: the agent maintains context across multiple steps.
  \item \textbf{Iteration}: the agent loops until the task is solved or a stopping condition is reached.
\end{itemize}

% ============================================================
\section{The ReAct pattern}

ReAct (Reasoning + Acting) interleaves reasoning traces with tool actions:

\begin{enumerate}
  \item \textbf{Thought}: the LLM reasons about the current state and next step.
  \item \textbf{Action}: the LLM calls a tool with specific inputs.
  \item \textbf{Observation}: the tool returns a result.
  \item Repeat until the LLM has enough information to answer.
\end{enumerate}

\begin{minted}{python}
# ReAct-style prompt (conceptual)
react_prompt = """Answer the question using the available tools.

Tools:
- search(query): search the web for information
- calculator(expression): evaluate a math expression

Question: What is the population of Benin multiplied by 3?

Thought: I need to find the population of Benin first.
Action: search("population of Benin 2025")
Observation: The population of Benin is approximately 14.4 million.
Thought: Now I need to multiply 14.4 million by 3.
Action: calculator("14400000 * 3")
Observation: 43200000
Thought: I have the answer.
Answer: The population of Benin multiplied by 3 is approximately 43,200,000."""

print(react_prompt)
\end{minted}

% ============================================================
\section{Function calling with Groq}

Modern LLM APIs support structured tool/function calling:

\begin{minted}{python}
from groq import Groq
import json

client = Groq()

# Define tools
tools = [
    {
        "type": "function",
        "function": {
            "name": "get_weather",
            "description": "Get the current weather for a city",
            "parameters": {
                "type": "object",
                "properties": {
                    "city": {"type": "string", "description": "City name"},
                    "unit": {"type": "string", "enum": ["celsius", "fahrenheit"]},
                },
                "required": ["city"],
            },
        },
    },
    {
        "type": "function",
        "function": {
            "name": "calculate",
            "description": "Evaluate a mathematical expression",
            "parameters": {
                "type": "object",
                "properties": {
                    "expression": {"type": "string", "description": "Math expression"},
                },
                "required": ["expression"],
            },
        },
    },
]

# Actual tool implementations
def get_weather(city, unit="celsius"):
    # In production, call a real weather API
    return {"city": city, "temperature": 28, "unit": unit, "condition": "sunny"}

def calculate(expression):
    return {"result": eval(expression)}  # caution: eval is unsafe in production

tool_map = {"get_weather": get_weather, "calculate": calculate}

# Agent loop
messages = [{"role": "user", "content": "What is the temperature in Cotonou in Fahrenheit?"}]

response = client.chat.completions.create(
    model="llama-3.1-8b-instant",
    messages=messages,
    tools=tools,
    tool_choice="auto",
)

# Process tool calls
msg = response.choices[0].message
if msg.tool_calls:
    messages.append(msg)
    for tc in msg.tool_calls:
        func_name = tc.function.name
        args = json.loads(tc.function.arguments)
        result = tool_map[func_name](**args)
        messages.append({
            "role": "tool",
            "tool_call_id": tc.id,
            "content": json.dumps(result),
        })
    # Get final response
    final = client.chat.completions.create(
        model="llama-3.1-8b-instant", messages=messages
    )
    print(final.choices[0].message.content)
\end{minted}

% ============================================================
\section{LangChain agents}

LangChain provides a higher-level agent framework:

\begin{minted}{python}
from langchain_google_genai import ChatGoogleGenerativeAI
from langchain_core.tools import tool
from langchain.agents import create_tool_calling_agent, AgentExecutor
from langchain_core.prompts import ChatPromptTemplate

# Define tools
@tool
def search_arxiv(query: str) -> str:
    """Search arxiv.org for recent papers. Returns titles and abstracts."""
    import urllib.request, json
    url = f"http://export.arxiv.org/api/query?search_query=all:{query}&max_results=3"
    response = urllib.request.urlopen(url).read().decode()
    # Simplified parsing
    return response[:2000]

@tool
def python_calculator(expression: str) -> str:
    """Evaluate a Python math expression. Example: '2**10 + 3*4'"""
    try:
        result = eval(expression)
        return str(result)
    except Exception as e:
        return f"Error: {e}"

# Create agent
llm = ChatGoogleGenerativeAI(model="gemini-2.0-flash", temperature=0)
tools_list = [search_arxiv, python_calculator]

prompt = ChatPromptTemplate.from_messages([
    ("system", "You are a helpful research assistant. Use tools when needed."),
    ("human", "{input}"),
    ("placeholder", "{agent_scratchpad}"),
])

agent = create_tool_calling_agent(llm, tools_list, prompt)
executor = AgentExecutor(agent=agent, tools=tools_list, verbose=True)

result = executor.invoke({"input": "Find recent papers about LoRA fine-tuning and tell me how many authors the first paper has."})
print(result["output"])
\end{minted}

% ============================================================
\section{Multi-step reasoning agents}

Agents can chain multiple tool calls to solve complex tasks:

\begin{minted}{python}
@tool
def read_file(path: str) -> str:
    """Read the contents of a text file."""
    with open(path) as f:
        return f.read()[:5000]

@tool
def summarize_text(text: str) -> str:
    """Summarize a long text into 3 bullet points."""
    llm = ChatGoogleGenerativeAI(model="gemini-2.0-flash", temperature=0.3)
    response = llm.invoke(f"Summarize in 3 bullet points:\n{text}")
    return response.content

@tool
def write_file(path: str, content: str) -> str:
    """Write content to a file."""
    with open(path, "w") as f:
        f.write(content)
    return f"Written to {path}"

tools_list = [read_file, summarize_text, write_file]
agent = create_tool_calling_agent(llm, tools_list, prompt)
executor = AgentExecutor(agent=agent, tools=tools_list, verbose=True)

result = executor.invoke({
    "input": "Read notes.txt, summarize it, and save the summary to summary.txt"
})
\end{minted}

% ============================================================
\section{LangGraph basics}

LangGraph enables building stateful, multi-step agent workflows as graphs:

\begin{minted}{python}
from langgraph.graph import StateGraph, END
from typing import TypedDict, Annotated

class AgentState(TypedDict):
    question: str
    research: str
    draft: str
    review: str
    final_answer: str

def research_step(state: AgentState) -> AgentState:
    """Research the question using the LLM."""
    llm = ChatGoogleGenerativeAI(model="gemini-2.0-flash")
    result = llm.invoke(f"Research this topic thoroughly: {state['question']}")
    return {"research": result.content}

def draft_step(state: AgentState) -> AgentState:
    """Write a draft answer based on research."""
    llm = ChatGoogleGenerativeAI(model="gemini-2.0-flash")
    result = llm.invoke(
        f"Based on this research:\n{state['research']}\n\n"
        f"Write a clear, concise answer to: {state['question']}"
    )
    return {"draft": result.content}

def review_step(state: AgentState) -> AgentState:
    """Review and improve the draft."""
    llm = ChatGoogleGenerativeAI(model="gemini-2.0-flash")
    result = llm.invoke(
        f"Review this draft for accuracy and clarity. Suggest improvements:\n"
        f"{state['draft']}"
    )
    return {"review": result.content}

def finalize_step(state: AgentState) -> AgentState:
    """Produce the final answer incorporating review feedback."""
    llm = ChatGoogleGenerativeAI(model="gemini-2.0-flash")
    result = llm.invoke(
        f"Original draft:\n{state['draft']}\n\n"
        f"Review feedback:\n{state['review']}\n\n"
        f"Write the final polished answer."
    )
    return {"final_answer": result.content}

# Build the graph
workflow = StateGraph(AgentState)
workflow.add_node("research", research_step)
workflow.add_node("draft", draft_step)
workflow.add_node("review", review_step)
workflow.add_node("finalize", finalize_step)

workflow.set_entry_point("research")
workflow.add_edge("research", "draft")
workflow.add_edge("draft", "review")
workflow.add_edge("review", "finalize")
workflow.add_edge("finalize", END)

app = workflow.compile()

result = app.invoke({"question": "What are the key differences between LoRA and QLoRA?"})
print(result["final_answer"])
\end{minted}

\begin{aitip}
LangGraph is more powerful than simple LangChain agents for workflows that need conditional branching, loops, or human-in-the-loop checkpoints. Use LangChain agents for simple tool use; use LangGraph for complex multi-step workflows.
\end{aitip}

\begin{exercise}
\begin{enumerate}
  \item Build a LangChain agent with 3 tools: web search, calculator, and current date/time. Test it with 5 queries that require different tool combinations.
  \item Implement function calling with the Groq API. Create a tool that looks up country information (capital, population, GDP) and an agent that answers geography questions.
  \item Build a LangGraph workflow with 4 steps: (1) research, (2) outline, (3) write, (4) review. Use it to generate a short article on any topic.
  \item Create a ``code assistant'' agent that can write Python code, execute it, and debug errors automatically.
  \item Compare the behavior of a simple LangChain agent vs a LangGraph workflow for the same complex task. Which produces better results?
\end{enumerate}
\end{exercise}

\begin{ethicsbox}
Agents can take real-world actions: send emails, modify files, execute code, call APIs. This amplifies both the benefits and the risks of AI. Implement these safeguards:
\begin{itemize}
  \item \textbf{Least privilege:} give agents only the tools they need.
  \item \textbf{Human approval:} require human confirmation for high-stakes actions (sending emails, deleting data).
  \item \textbf{Sandboxing:} run code execution in isolated environments.
  \item \textbf{Rate limiting:} prevent agents from making unlimited API calls or taking unlimited actions.
  \item \textbf{Audit logging:} log every action the agent takes for review.
\end{itemize}
\end{ethicsbox}

% ============================================================
\section{Chapter summary}

\begin{itemize}
  \item An LLM agent combines reasoning, tool use, memory, and iteration.
  \item The ReAct pattern interleaves thinking and acting in a loop.
  \item Function calling lets LLMs invoke structured tools via API.
  \item LangChain agents provide a high-level framework for tool-using LLMs.
  \item LangGraph enables stateful, multi-step workflows with conditional branching.
  \item Agent safety requires least privilege, human oversight, sandboxing, and audit logging.
\end{itemize}
